% 第3章 研究内容、技术路线、创新点

\chapter{三、研究内容、技术路线、创新点}

\section{（1）主要研究内容}

\subsection{1）中红外通道吸收池光路联调与浓度反演精度验证}

围绕四通道中红外 QCL-TDLAS 平台，拟重点突破折返式吸收池光路（原位反射框架）的装调与``吸收谱—温度—浓度''联合反演的精度验证：（1）完成反射镜组排布、光轴对准与有效光程标定，解决窗口污染引起的透过率漂移补偿与光束抖动对反演稳定性的影响；（2）基于 Voigt 线型前向模型，由双线线强比或 H$_2$O 通道确定温度、以非线性最小二乘反演各组分浓度，系统评估谱线强度、配分函数、压力展宽系数、基线漂移与光程标定误差对反演结果的传递规律，给出各通道的浓度不确定度预算；（3）结合 Cantera 一维火焰仿真完成浓度/温度标定，并在不同当量比、不同停留时间下考察模型的外推能力与高温高压工况下的反演精度（目标：实验—仿真偏差 $<$ 20\%）。

\subsection{2）探明 QCL-TDLAS 与 NO-PLIF 两种手段在发动机尾气工况下的耦合与互校验机制，形成一体化诊断能力}

针对多组分 TDLAS 与自由基 PLIF 长期各自独立、缺乏同一标准场下交叉校验的问题，拟建立两者的耦合与互校验机制：

（1）以统一主时钟/触发源同步 QCL 波长扫描时序、burst 染料激光脉冲串与高速相机曝光时序，实现``同一时刻、同一区域''的同步测量；

（2）利用多程吸收光路与片光的几何对应关系，把 TDLAS 路径积分浓度映射到 PLIF 二维浓度场，由 PLIF 提供空间均匀性先验以校正``路径平均浓度''与``峰值浓度''差异，反向以 TDLAS 免标定结果为 PLIF 强度—浓度标定提供基准，形成双向互校验；

（3）在标准火焰与发动机尾气工况下建立互校验模型与合成不确定度框架，最终形成``线定量 + 场成像''相互印证的一体化 NO 诊断能力。

\section{拟采取的研究方法、技术路线、实施方案及可行性分析}

\subsection{（1）研究方法}

本课题采用``仿真选线—平台搭建—标准场标定—现场验证''的研究方法：首先以 HITRAN/HITEMP 仿真与 Cantera 一维火焰计算确定谱线与工况；其次并行搭建中红外 QCL-TDLAS 与 burst 染料激光 PLIF 两套平台；随后在标准火焰台完成标定；最后在地面试验台/发动机尾气工况下开展现场测量与交叉校验，并对实验—仿真偏差进行校正。

\subsection{（2）技术路线与实施方案}

\subsubsection{1）技术路线}

本课题采用``仿真选线—平台搭建—标准场标定—现场验证''的技术路线：首先以 HITRAN/HITEMP 仿真与 Cantera 一维火焰计算确定谱线与工况；其次并行搭建中红外 QCL-TDLAS 与 burst 染料激光 PLIF 两套平台；随后在 Hencken 标准火焰台完成标定；最后在地面试验台/发动机尾气工况下开展现场测量与交叉校验，并对实验—仿真偏差进行校正。

\subsubsection{2）实施方案}

阶段一：仿真与选线。基于 HITRAN/HITEMP 与自编 Voigt 程序筛选 NO 目标谱线，用 Cantera（GRI-Mech 3.0）一维 Hencken 火焰计算目标工况的温度/物种分布，为选线与反演提供先验。

阶段二：TDLAS 平台搭建（结构设计 + 联调）。完成 QCL 发射—探测、折返式长光程吸收池（原位反射框架）装调与有效光程标定、波长扫描与联合反演算法，结合 Cantera 完成浓度/温度标定与误差评估。

阶段三：PLIF 平台搭建（结构设计 + 联调）。完成染料激光泵浦—振荡光路、高压染料循环/蓄能系统（染料循环泵系统制作）、283 nm 片光与高速相机采集；攻克 burst 模式能量稳定性、线宽压窄与频率锁定，完成强度—浓度标定。

阶段四：点—面协同与交叉校验。统一时序触发，在标准火焰与发动机尾气工况下开展 TDLAS 与 PLIF 交叉比对、偏差校正与不确定度合成，形成一体化诊断能力。

\subsection{（3）可行性分析}

\subsubsection{1）理论和技术可行}

以 Voigt 线型为前向模型的``吸收谱—温度—浓度''联合反演、PLIF 荧光猝灭校正、标准火焰标定等均有成熟物理基础与大量文献支撑；burst 染料激光与 QCL 技术已工程化；本课题提出的``线定量 + 场成像''互校验架构在物理与信息学上自洽，理论可行。中红外 QCL 四通道激光器与 burst 染料激光平台均已到位，关键技术路线清晰。

\subsubsection{2）实验条件可行}

课题组已具备：中红外 QCL 通道激光器（NO）与 VIGO 中红外探测器（已验收）；HITRAN/HITEMP + 自编 Voigt 仿真代码、Cantera 一维火焰仿真；染料泵系统、高速相机、片光源等已到货；既有激光诊断平台。关键仪器与经费基本到位，实验条件可行。

\section{论文的创新点}

（1）提出``中红外多组分 TDLAS 线测量 + burst 染料激光自由基面测量''的协同诊断架构，实现同一燃烧场内浓度与结构的同步表征。

（2）针对 burst 模式下染料激光能量抖动与线宽展宽问题，发展高压染料循环稳流与蓄能缓冲相结合的稳定性控制方法（染料循环泵系统与蓄能缓冲设计）。

（3）建立多组分吸收谱线与 PLIF 自由基图像在标准火焰场下的联合标定与交叉校验方法，形成``线定量 + 场成像''相互印证的一体化 NO 诊断能力——即探明 QCL-TDLAS 与 NO-PLIF 在发动机尾气工况下的耦合与互校验机制。
